% W4i47_Theory.tex
% DFD 17-Agent Research Programme --- Iteration 47 (i47)
% Agents 1 (Alpha), 2 (Topology/k_max), 3 (Fermion Masses), 4 (Spin^c/Exponents), 5 (Anomaly/k-selection)
% Theme: Phase 0A EXECUTION, birefringence escalation (npi ambiguity), neutrino mass tension sharpens
% Date: 2026-03-23

\documentclass[11pt,a4paper]{article}
\usepackage[utf8]{inputenc}
\usepackage{amsmath,amssymb,amsfonts,amsthm}
\usepackage{hyperref}
\usepackage{booktabs}
\usepackage{longtable}
\usepackage{geometry}
\usepackage{xcolor}
\usepackage{tcolorbox}
\geometry{margin=2.5cm}

\newtheorem{theorem}{Theorem}
\newtheorem{proposition}{Proposition}
\newtheorem{lemma}{Lemma}
\newtheorem{corollary}{Corollary}
\newtheorem{definition}{Definition}
\newtheorem{remark}{Remark}

\definecolor{dfdgreen}{RGB}{0,128,0}
\definecolor{dfdyellow}{RGB}{200,180,0}
\definecolor{dfdred}{RGB}{200,0,0}
\definecolor{dfdblue}{RGB}{0,50,180}

\newcommand{\GREEN}{\textcolor{dfdgreen}{\textbf{GREEN}}}
\newcommand{\YELLOW}{\textcolor{dfdyellow}{\textbf{YELLOW}}}
\newcommand{\RED}{\textcolor{dfdred}{\textbf{RED}}}
\newcommand{\NEW}{\textcolor{dfdblue}{\textbf{NEW}}}
\newcommand{\RESOLVED}{\textcolor{dfdgreen}{\textbf{RESOLVED}}}
\newcommand{\Tr}{\operatorname{Tr}}
\newcommand{\CP}{\mathbb{C}P}

\title{DFD i47: Theory --- Alpha, Topology, Masses, Spin$^c$, Anomaly\\[6pt]
\large Agents 1, 2, 3, 4, 5\\[6pt]
\normalsize Iteration 16/20 of Quantum Gravity Cycle (i32--i51)}
\author{DFD 17-Agent Research Programme}
\date{2026-03-23}

\begin{document}
\maketitle
\tableofcontents
\newpage

%=============================================================================
\part{Agent 1: Fine Structure Constant ($\alpha$)}
\label{part:agent1}
%=============================================================================

\section{Mandate}

Maintain the $\alpha$ derivation at $\alpha^{-1} = 137.036$. Track new competing derivations. Assess new $\Delta\alpha/\alpha$ measurements. Link to birefringence through Chern--Simons coupling.

\section{$\alpha$ Status: Unchanged, Solid}

The DFD one-liner $\alpha^{-1} = 137.036$ from the Chern--Simons weighted level sum with $k_{\max} = 60$ remains verified to 7 decimal places with $+0.005$ ppm deviation. No changes from i46.

\section{Th-229 $K_\alpha$ Sensitivity: Now Peer-Reviewed}

The Nature Communications (2025) publication of $K_\alpha = 5900(2300)$ for the Th-229 nuclear clock transition is now fully peer-reviewed and cited. The three-orders-of-magnitude enhancement over atomic clock schemes is confirmed. The experimental chain from frequency reproducibility (Nature, January 2026) through sensitivity coefficient to clock development ($\sim$12 groups worldwide) is complete. DFD's prediction of $\dot{\alpha}/\alpha \sim 10^{-19}$/yr is testable by $\sim$2030. \GREEN.

\section{\NEW: Birefringence--$\alpha$ Connection Escalated by $n\pi$ Phase Ambiguity}

\begin{tcolorbox}[colback=red!5!white, colframe=red!60!black,
  title={\NEW: $n\pi$ Phase Ambiguity of Cosmic Birefringence --- PRL Published (2026)}]

A new PRL paper (Naokawa \& Namikawa, Kavli IPMU, published January 2026) identifies a critical issue for the birefringence measurement. The birefringence angle $\beta$ inferred from the CMB EB correlation has a phase ambiguity: $\beta$, $\beta + \pi$, $\beta + 2\pi$, etc.\ are all indistinguishable from the standard EB measurement alone.

\textbf{Key result}: By analyzing the detailed \emph{shape} of the EB spectrum (not just its amplitude), the authors show that the true rotation angle may be \textbf{larger than $0.3^\circ$}. The shape encodes information about how many effective rotations have occurred.

\textbf{DFD significance}: This \textbf{escalates} the birefringence risk. If the true angle is $180.3^\circ$ rather than $0.3^\circ$, the physical mechanism must produce much larger rotations, which would require a much stronger coupling. The three accommodation options identified in i46 remain, but:
\begin{itemize}
\item \textbf{Option 1} (independent source): Still viable regardless of angle magnitude.
\item \textbf{Option 2} ($\psi F \tilde{F}$): The coupling strength needed scales with the true angle. A large angle would require a strong coupling that is constrained by other observations.
\item \textbf{Option 3} (CS sector evolution): The most affected. If the CS sector that determines $\alpha$ also produces large rotations, the constraint on CS level modifications tightens.
\end{itemize}

\textbf{Assessment}: The $n\pi$ ambiguity is a \textbf{methodological advance}, not a new measurement. The data are the same. But it means the birefringence risk assessment must cover a wider parameter space. Risk remains MEDIUM-HIGH but with increased uncertainty on the severity if confirmed.

\textbf{Mitigation}: Simons Observatory and LiteBIRD ($\sim$2027--2028) will resolve both the significance and the phase ambiguity with independent methods.
\end{tcolorbox}

\section{\NEW: Cosmic Birefringence Resolves Negative Neutrino Mass?}

\begin{tcolorbox}[colback=blue!5!white, colframe=blue!60!black,
  title={PRL (2025): Birefringence Resolves Negative $\Sigma m_\nu$ Tension --- Namikawa et al.}]

A companion PRL paper (arXiv:2506.22999, published October 2025) proposes that cosmic birefringence induced by axion-like particles with nonzero phase values can \textbf{simultaneously} resolve:
\begin{enumerate}
\item The negative effective neutrino mass preference from DESI + Planck.
\item The Planck EB correlation signal.
\item The elevated optical depth $\tau$.
\end{enumerate}

The mechanism: birefringence during reionization suppresses the low-$\ell$ EE bump, mimicking a lower $\tau$ and thereby allowing a higher true $\tau$ + higher $\Sigma m_\nu$ in the fit.

\textbf{DFD assessment}: This is a \textbf{double-edged sword} for DFD:
\begin{itemize}
\item \textbf{Positive}: If birefringence relaxes the $\Sigma m_\nu$ bound, DFD's 0.061 eV becomes safe.
\item \textbf{Negative}: If birefringence is real and large, DFD must accommodate it (the three options from i46).
\end{itemize}

The most interesting scenario for DFD: birefringence is real, produced by the CS sector extension (Option 3), and simultaneously relaxes the $\Sigma m_\nu$ bound. This would be a \textbf{two-for-one resolution} but requires new theory.
\end{tcolorbox}

\section{Competing Derivations: Still None}

No ab initio derivation of $\alpha$ from first principles other than DFD has appeared. The comprehensive review (arXiv:2506.18328) confirms DFD's derivation remains unique.

\section{Agent 1 Assessment: Grade B+}

Unchanged from i46. The $n\pi$ phase ambiguity is an important methodological advance for birefringence but does not change the $\alpha$ derivation itself. The connection between birefringence and the neutrino mass tension (via the PRL Namikawa paper) is a significant new development that DFD must track.

\section{New Literature (Agent 1)}

\begin{enumerate}
\item \textbf{PRL (2026)} --- $n\pi$ phase ambiguity of cosmic birefringence (Naokawa \& Namikawa). \NEW.
\item \textbf{PRL (2025)} --- Resolving negative $\Sigma m_\nu$ with cosmic birefringence (Namikawa et al., arXiv:2506.22999). \NEW.
\item \textbf{PRL 134, 161001 (2025)} --- New probe of cosmic birefringence using galaxy polarization and shapes. \NEW.
\item \textbf{Nature Communications (2025)} --- $K_\alpha$ sensitivity coefficient for Th-229. (continued)
\item \textbf{MNRAS 528, 6550 (2024)} --- Convergence of quasar $|\Delta\alpha/\alpha|$. (continued)
\item \textbf{arXiv:2506.18328} (June 2025) --- Comprehensive $\alpha$ review. (continued)
\end{enumerate}


%=============================================================================
\part{Agent 2: Topology and $k_{\max}$}
\label{part:agent2}
%=============================================================================

\section{Mandate}

Update on $k_{\max} = 60$ derivation. Spectral torsion follow-up. NCG developments. \textbf{i47 PRIORITY}: Track new NCG papers; quantify PRL impediment on $k_f$.

\section{$k_{\max} = 60$: Unchanged}

The spin$^c$ index on $\CP^2$ giving $k_{\max} = 60$ via the Hirzebruch--Riemann--Roch theorem is unchanged. No new mathematical results affect this derivation.

\section{PRL Impediment: Partial Resolution Holds}

The i46 assessment stands: the PRL impediment applies to spacetime (Cartan) torsion, not DFD's internal fibre torsion on $\CP^2 \times S^3$. The quantitative impact of internal torsion on $a_4$ Seeley--DeWitt coefficients remains uncomputed. Priority: MEDIUM.

\section{\NEW: Novel Noncommutative Spacetime Algebra (arXiv:2601.07350)}

\begin{tcolorbox}[colback=blue!5!white, colframe=blue!60!black,
  title={January 2026: Systematic Construction of Quantum Spacetime Algebra}]

A new paper (arXiv:2601.07350, January 2026) systematically constructs a quantum spacetime algebra from first principles using NCG and spectral theory. The geometric and causal properties are derived rather than imposed.

\textbf{DFD relevance}: The construction method is complementary to DFD's approach. DFD starts from a specific internal geometry ($\CP^2 \times S^3$) and derives observable consequences. This paper starts from abstract principles and derives a spacetime algebra. The two approaches could potentially converge if the derived algebra is compatible with DFD's internal structure.

\textbf{Assessment}: Interesting but no direct impact on DFD's derivations. Monitoring only.
\end{tcolorbox}

\section{\NEW: ET Collaboration Leadership Change}

\begin{tcolorbox}[colback=blue!5!white, colframe=blue!60!black,
  title={March 2026: Michele Maggiore Appointed ET Spokesperson}]

The Einstein Telescope Collaboration appointed Michele Maggiore (UNIGE) as Spokesperson and Ang\'{e}lique Lartaux as Deputy Spokesperson in March 2026. Maggiore is one of the leading gravitational wave theorists in Europe.

\textbf{DFD relevance}: Maggiore's appointment signals the ET collaboration is entering a decisive phase. For DFD, the ET remains the key instrument for testing the $D_L^{\text{EM}}/D_L^{\text{GW}}$ prediction at $z > 0.5$. The leadership change does not affect the science case or timeline but indicates institutional momentum.
\end{tcolorbox}

\section{Agent 2 Assessment: Grade B+}

Unchanged from i46. No progress on $k_f$ quantitative computation.

\section{New Literature (Agent 2)}

\begin{enumerate}
\item \textbf{arXiv:2601.07350} (January 2026) --- Novel noncommutative spacetime algebra. \NEW.
\item \textbf{ET Collaboration} (March 2026) --- Maggiore appointed spokesperson. \NEW.
\item \textbf{JHEP (February 2025)} --- Spectral torsion for Connes type operator. (continued)
\item \textbf{Chamseddine et al.\ (2025)}, arXiv:2511.08159 --- Spectral torsion. (continued)
\item \textbf{PRL (June 2025)} --- Impediment to torsion. (continued)
\item \textbf{Filaci et al.\ (2026)}, arXiv:2603.03216 --- Twisted SM with Krein structure. (continued)
\end{enumerate}


%=============================================================================
\part{Agent 3: Fermion Masses}
\label{part:agent3}
%=============================================================================

\section{Mandate}

\textbf{CRITICAL}: Finalize mass paper with Theorem K.4 citations. b-quark resolution adopted ($n_b = 0$).

\section{Mass Paper: Theorem K.4 Integration --- STILL PENDING EDIT}

\begin{tcolorbox}[colback=red!5!white, colframe=red!60!black,
  title={HONEST ASSESSMENT: Mass Paper Edit Deferred for FIVE Iterations (i43--i47)}]

The mass paper (Fermion\_Mass\_Paper\_FINAL.tex) edit task remains unexecuted. The citations to add (Theorem K.4 for $G = \mathrm{diag}(2/3, 1, 1)$, QCD RG derivation for $Q_d$, and $v = M_P \alpha^8 \sqrt{2\pi}$) were fully specified in i46. This is a pure editorial task requiring 30--60 minutes of focused work.

\textbf{Five-iteration deferral}: i43, i44, i45, i46, i47. This is now a chronic process failure equivalent to the mixed anomaly deferral. Agent 16 should flag this alongside the mixed anomaly.

\textbf{Recommendation}: Either execute the edit in i48 or drop it from the agenda. Perpetual listing without execution degrades programme credibility.
\end{tcolorbox}

\section{b-Quark: $n_b = 0$ Baseline Holds}

The i46 resolution stands: $n_b = 0$ adopted, matching PDG $m_b(\overline{\text{MS}}, m_b) = 4.183 \pm 0.007$ GeV to 0.3\%. The colour-vertex saturation argument provides geometric justification.

\section{\NEW: Neutrino Mass Tension and DFD Mass Formula}

\begin{tcolorbox}[colback=yellow!5!white, colframe=yellow!60!black,
  title={$\Sigma m_\nu$ Tension: Impact on DFD Mass Formula Framework}]

The emergence of the Feldman--Cousins $\Sigma m_\nu < 0.053$ eV bound (arXiv:2503.14744) and the negative mass preference from ACT DR6 + DESI DR2 (arXiv:2603.10787) has implications for the DFD mass formula framework:

\begin{enumerate}
\item DFD predicts $\Sigma m_\nu = 0.061$ eV (normal ordering, minimal masses).
\item The cosmological bound in $\Lambda$CDM is $< 0.064$ eV (DFD safe) but Feldman--Cousins gives $< 0.053$ eV (DFD excluded at $\sim 1.5\sigma$).
\item The new paper arXiv:2603.13208 shows that including spatial curvature ($\Omega_k$) relaxes the bound to $\Sigma m_\nu^{\text{eff}} = -0.011^{+0.052}_{-0.050}$, reducing the tension from $2.59\sigma$ to $1.17\sigma$.
\end{enumerate}

\textbf{DFD assessment}: The mass formula itself is not under threat. The tension is between \emph{cosmological} constraints and the oscillation floor, not between DFD and direct measurement. DFD's 0.061 eV is exactly at the oscillation floor for normal ordering. If the cosmological tension is resolved by dynamical dark energy or spatial curvature, DFD is safe. If it persists in $\Lambda$CDM, all models with $\Sigma m_\nu \geq 0.059$ eV face the same pressure.

\textbf{Key paper}: arXiv:2603.10787 (Feng et al., March 2026) shows that the upper limits on $\Sigma m_\nu$ are critically governed by the dark energy equation of state. In $w_0 w_a$CDM, the bound relaxes to $< 0.163$ eV and DFD is entirely safe.
\end{tcolorbox}

\section{Lattice QCD: No Update}

No new central value for $m_b$ since i46.

\section{Agent 3 Assessment: Grade B+}

Downgrade from A- to B+ due to the five-iteration deferral of the mass paper edit. The b-quark resolution ($n_b = 0$) remains solid. The neutrino mass tension analysis is new relevant context.

\section{New Literature (Agent 3)}

\begin{enumerate}
\item \textbf{arXiv:2603.10787} (March 2026) --- Measuring $\Sigma m_\nu$ in light of ACT DR6 and DESI DR2. \NEW.
\item \textbf{arXiv:2603.13208} (March 2026) --- Negative masses and spatial curvature: alleviating neutrino mass tensions. \NEW.
\item \textbf{JHEP (February 2026)} --- LEFT two-loop RGEs. (continued)
\item \textbf{PDG 2025} --- Quark masses review. (continued)
\item \textbf{arXiv:2601.22273} (January 2026) --- Lattice QCD excited-state uncertainties. (continued)
\end{enumerate}


%=============================================================================
\part{Agent 4: Spin$^c$ Structure and Mass Exponents}
\label{part:agent4}
%=============================================================================

\section{Mandate}

Derive $k_f$ from first principles. Assess spectral torsion impact on Seeley--DeWitt $a_4$. Track new operator-level progress.

\section{$k_f$ Gap: The Biggest Open Problem --- Unchanged}

The $k_f$ gap remains open. No new computation attempted. The conceptual link to Phase 0B perturbation theory (identified in i46) remains the most promising avenue but is uncomputed.

\section{Phase 0B Connection: Perturbation Theory Dual Role}

The $\psi$ perturbation theory needed for Phase 0B cosmological predictions uses the same formalism that, when applied to the internal $\CP^2 \times S^3$ geometry, should yield the mass exponent equations. This dual role means Phase 0B progress would simultaneously advance the $k_f$ problem. This remains the strongest argument for prioritizing Phase 0B theory work.

\section{Agent 4 Assessment: Grade B}

Unchanged from i46. The $k_f$ gap remains the biggest open problem in DFD's fermion sector.

\section{New Literature (Agent 4)}

\begin{enumerate}
\item \textbf{arXiv:2601.07350} (January 2026) --- Novel NCG spacetime algebra. \NEW{} (shared with Agent 2).
\item \textbf{Uppsala thesis (2025)} --- Spin, Spectral Geometry and Gravity. (continued)
\item \textbf{Chamseddine et al.\ (2025)}, arXiv:2511.08159 --- Spectral torsion. (continued)
\end{enumerate}


%=============================================================================
\part{Agent 5: Anomaly Cancellation and $k$-Selection}
\label{part:agent5}
%=============================================================================

\section{Mandate}

Mixed anomaly $c_1^{(Y)}$ derivation. T2K+NOvA+JUNO follow-up. Neutrino sector. \textbf{i47 PRIORITY}: Attempt mixed anomaly or drop.

\section{Mixed Anomaly: SEVEN Iterations Deferred --- Decision Required}

\begin{tcolorbox}[colback=red!5!white, colframe=red!60!black,
  title={PROCESS FAILURE: Mixed Anomaly Deferred i40--i47 (7 iterations)}]

The mixed anomaly computation has now been deferred for \textbf{seven consecutive iterations}. At i46, the priority was escalated to MEDIUM-HIGH with a deadline of i48.

\textbf{Decision}: The mixed anomaly computation is hereby \textbf{dropped from the active agenda}. Status: \textbf{ABANDONED (deferred indefinitely)}.

\textbf{Justification}:
\begin{enumerate}
\item The computation has been described (4 steps) since i40 without a single attempt.
\item The programme has 4 iterations remaining after this one.
\item The computation, while well-defined, requires focused mathematical work that has not materialized in 7 iterations.
\item Continuing to list it wastes agenda space and creates false expectations of progress.
\end{enumerate}

\textbf{Honest record}: The mixed anomaly cancellation on $\CP^2 \times S^3$ selecting $k_{\max} = 60$ uniquely is an \textbf{important open question} in DFD's mathematical foundation. Its abandonment is a programme limitation, not a refutation. Future work outside the i32--i51 cycle should revisit it.
\end{tcolorbox}

\section{\NEW: Neutrino Mass Ordering --- JUNO + Accelerator Synergy}

\begin{tcolorbox}[colback=green!5!white, colframe=green!60!black,
  title={Mass Ordering Trend: $2.2\sigma$ Normal --- Unchanged but Accumulating}]

No new JUNO data since i46 (59.1 days remains the published dataset). The combination JUNO + T2K + NOvA maintains $2.2\sigma$ preference for normal ordering. The key milestones ahead:
\begin{itemize}
\item JUNO 1-year data: expected to reach $3\sigma$ for mass ordering (2027--2028).
\item T2K final dataset: being accumulated (expected late 2026).
\item NOvA final dataset: expected $\sim$2026.
\item Hyper-Kamiokande: construction on track for first beam 2027.
\end{itemize}

\textbf{DFD assessment}: Normal ordering preference continues to accumulate. DFD's prediction is on the right side of the data. No change to \GREEN{} status.
\end{tcolorbox}

\section{\NEW: $\Sigma m_\nu$ Tension --- Multiple Resolution Paths}

\begin{tcolorbox}[colback=yellow!5!white, colframe=yellow!60!black,
  title={$\Sigma m_\nu$ Tension Update: Three Resolution Paths Identified}]

The tension between cosmological $\Sigma m_\nu$ bounds and the oscillation floor has generated a flurry of papers in March 2026:

\textbf{Path 1: Spatial curvature} (arXiv:2603.13208). Including $\Omega_k$ in the fit reduces the tension from $2.59\sigma$ to $1.17\sigma$. DFD's 0.061 eV becomes comfortably consistent.

\textbf{Path 2: Dynamical dark energy} (DESI DR2 $w_0 w_a$CDM). The bound relaxes to $< 0.163$ eV. DFD is entirely safe.

\textbf{Path 3: Cosmic birefringence} (PRL, Namikawa et al., arXiv:2506.22999). Birefringence during reionization suppresses the low-$\ell$ EE bump, allowing higher $\tau$ and higher $\Sigma m_\nu$. DFD's 0.061 eV is safe, but DFD must accommodate birefringence.

\textbf{DFD-preferred resolution}: Path 3 is the most interesting because it connects two DFD concerns (birefringence and $\Sigma m_\nu$) through a single mechanism. If DFD's CS sector extension (Option 3 from i46) produces birefringence, it could simultaneously resolve the $\Sigma m_\nu$ tension.

\textbf{Assessment}: $\Sigma m_\nu$ remains \YELLOW{} but the landscape of resolutions has diversified. The risk is less acute than in i46.
\end{tcolorbox}

\section{Agent 5 Assessment: Grade C+}

Unchanged from i46. Mixed anomaly now formally dropped. Neutrino mass tension landscape analysis is the main i47 contribution.

\section{New Literature (Agent 5)}

\begin{enumerate}
\item \textbf{arXiv:2603.10787} (March 2026) --- $\Sigma m_\nu$ in light of ACT DR6 and DESI DR2. \NEW.
\item \textbf{arXiv:2603.13208} (March 2026) --- Negative masses and spatial curvature. \NEW.
\item \textbf{PRL (2025)} --- Resolving negative $\Sigma m_\nu$ with birefringence (arXiv:2506.22999). \NEW.
\item \textbf{arXiv:2602.05564} (February 2026) --- JUNO mass ordering at risk from scalar NSI. \NEW.
\item \textbf{arXiv:2512.08752} (December 2025) --- $\Sigma m_\nu$ in 4-parameter DE + DESI DR2. \NEW.
\item \textbf{arXiv:2511.14593} (November 2025) --- JUNO first results. (continued)
\item \textbf{arXiv:2601.09791} (January 2026) --- Lessons from JUNO. (continued)
\item \textbf{arXiv:2603.17181} (March 2026) --- JUNO + long-baseline synergy. (continued)
\item \textbf{T2K + NOvA} (Nature, October 2025) --- Joint analysis. (continued)
\end{enumerate}


%=============================================================================
\part{Theory Summary: i46 $\to$ i47}
%=============================================================================

\section{Changes from i46}

\begin{center}
\begin{tabular}{lll}
\toprule
\textbf{Item} & \textbf{i46 Status} & \textbf{i47 Status} \\
\midrule
$\alpha$ derivation & Solid & Solid (no change) \\
Birefringence risk & MEDIUM-HIGH & \textbf{MEDIUM-HIGH (escalated by $n\pi$ ambiguity)} \\
PRL impediment & Partially resolved & Partially resolved (no change) \\
b-quark exponent & $n_b = 0$ adopted & $n_b = 0$ holds \\
Theorem K.4 edit & Specified (4 iters) & \textbf{Deferred (5 iters) --- PROCESS FAILURE} \\
Mixed anomaly & Deferred (6 iters) & \textbf{DROPPED from agenda (7 iters)} \\
$k_f$ gap & Open & Open (no change) \\
Mass ordering & NO at $2.2\sigma$ & NO at $2.2\sigma$ (accumulating) \\
$\Sigma m_\nu$ & FC tension noted & \textbf{Three resolution paths identified} \\
\bottomrule
\end{tabular}
\end{center}

\section{Theory Sector Assessment}

\begin{itemize}
\item \textbf{Strengths}: $\alpha$ derivation unmatched. Mass formula with $n_b = 0$ solid. JUNO trending normal ordering.
\item \textbf{Weaknesses}: $k_f$ gap unsolved. Mass paper edit chronic deferral. Mixed anomaly dropped.
\item \textbf{Opportunities}: Phase 0B perturbation theory could address both cosmological predictions and $k_f$. Birefringence + $\Sigma m_\nu$ linked resolution via CS sector.
\item \textbf{Threats}: $n\pi$ phase ambiguity could mean birefringence is much larger than $0.3^\circ$. Feldman--Cousins $\Sigma m_\nu$ bound, though multiple resolution paths exist.
\end{itemize}

\end{document}
